% ============================================================
% CORE MECHANICS — SECOND DRAFT (reconciled to character sheet)
% ------------------------------------------------------------
% RESOLVED BY THE SHEET (previous proposals corrected):
%   - Ratings run 0–10, not 0–5. Thresholds: 5+ = supernatural
%     influence; 8+ = masterful. (Species traits that grant "a
%     dot above the cap" now read against creation/tier caps,
%     TBD in the overhaul.)
%   - Success is FIXED at 4+ on a d6. The GM's difficulty knob is
%     REQUIRED SUCCESSES (and situational dice), not a moving
%     per-die threshold.
%   - Skill list is the sheet's 12, in 4 groups: Mobility
%     (Athletics, Knack, Nimbleness), Exploration (Investigate,
%     Perceive, Survive), Knowledge (Study, Folklore, Religion),
%     Social (Convince, Deceive, Intimidate).
%   - Combat: Strike (melee dice) / Shoot (ranged dice) / Defend
%     (SUBTRACTS dice from attacker) / Armor / Move / HP.
%   - Alignments have a second use: "may add dice when
%     applicable" (augment), alongside the spec's invocation.
%   - Identity fields: Demeanor, Drive, Reputation, Nickname/
%     Wanted-Poster Name. Deeds = 6 slots of specializations &
%     signature abilities. Possessions track Coins and Funds.
%
% REMAINING [PROPOSALS] FOR AUTHOR SIGN-OFF:
%   P1. Augment vs Invoke pricing: augmenting a Skill with
%       Alignment dice invites the plane (GM MAY complicate);
%       invoking pure Alignment (auto-successes = rating) ALWAYS
%       costs. Two uses, one escalating price.
%   P2. Saves: roll dice = rating in the targeted Alignment,
%       4+ succeeds, vs the effect's required successes; forced,
%       so no complication.
%   P3. Unspent Momentum lost at end of activation.
%   P4. Reload to 5 at the start of your side's activation.
%   P5. Damage scaffold: net hits (after Defend subtraction) =
%       damage; Armor reduces damage taken. Placeholder pending
%       the combat chapter.
%   P6. "Track" (Squar trait) maps to Survive; note added.
%   P7. Leveling: +2 cards/level from pools at or below your
%       level, 1 free swap. Deck grows; pruning intended.
%   P8. Unaffiliated: Drifter pool replaces faction cards.
%   P9. End of encounter: all cards incl. One-Shots re-holster.
%   P10. Coins = Rings in hand; Funds = banked credit and scrip
%        (welds to the Trade Companies economy). Recommend
%        unifying "Heaven"/"Heavens" and "Wyrd" spellings between
%        sheet and cosmology docs.
%   P11. Creation arrays (skill/alignment points): PLACEHOLDER.
%   P12. Deeds: earned at milestones, named for accomplished
%        feats; each grants a specialization die or signature
%        use. Sketched, not final.
% ============================================================

\chapter{How the Game Works}

\section{The Shape of Play}

You play one of the frontier's rare capable people --- the sort the factions recruit, the factors hire, and the ground remembers --- and the sheet measures you in three currencies.

\textbf{Skills} are what you have learned: twelve trained competences, rated 0--10 and rolled in dice. \textbf{Alignments} are what you can call on: the five great natures of the cosmos, expressed through you --- able to sharpen your skills when their nature applies, able to act outright when you ask them to, and never, in the second case, free. \textbf{Powers} are what you carry: a deck of cards built from your class, your species, and your faction, played from a loaded hand like rounds from a cylinder.

Everything in an encounter runs on \textbf{Momentum} --- a shared, swingy pool of action your whole side spends together --- and the rhythm of play is simple: your side surges, spends, and passes; theirs answers; and the fight is decided as much by how each side spends its luck as by what any one character can do.

One thing to know before the numbers: none of this machinery is arbitrary. The five Alignments are the five planes --- the Heavens, the Machine, the Shadow, the Wild, and the Wyrd --- and their pentacle on your sheet is a map of the actual cosmos. When you mark a dot in one, you are recording a relationship. When you use one, you are briefly conducting a power that holds the world's spinning top upright. The top is stable only in motion. So are you. That is what Momentum is.

\medskip\hrule\medskip

\section{Dice and Ratings}

The game uses six-sided dice, rolled in pools, and one rule underneath everything:

\begin{center}
\textbf{A die succeeds on a 4, 5, or 6. Count your successes.}
\end{center}

The target never moves. What moves is how many successes a task requires, and how many dice you bring.

\textbf{Ratings.} Skills and Alignments are rated \textbf{0--10}, and the scale means something in the world:

\begin{center}
\begin{tabular}{lll}
\textbf{Rating} & \textbf{Standing} & \\
\hline
1--2 & Capable & An ordinary practitioner \\
3--4 & Professional & The person the town sends for \\
\textbf{5--7} & \textbf{Supernatural influence} & Past what flesh and practice explain \\
\textbf{8--10} & \textbf{Masterful} & The handful per generation \\
\end{tabular}
\end{center}

The thresholds are not flavor. At 5+, a rating has crossed into territory the setting treats as supernatural --- the Guild examiner's impossible eye, the duelist the Corvara would call quick, the tracker the network vouches without a second's pause --- and effects, cards, and traits will reference the thresholds by name. At 8+, you are a fact of the region.

\textbf{Task resolution.} Roll dice equal to your Skill; count 4+ successes; compare to what the GM asked for:

\begin{center}
\begin{tabular}{ll}
\textbf{Required successes} & \textbf{Feels like} \\
\hline
1 & A real task with real stakes \\
2 & Hard; the frontier's usual answer \\
3 & A professional's bad afternoon \\
4--5 & Stories get told \\
\end{tabular}
\end{center}

Extra successes beyond the requirement are \textbf{margin}: faster, cleaner, quieter, \textit{more}. The GM converts margin into advantage, and players should feel free to propose what their margin buys. Situational edges and burdens (good tools, darkness, a Serekh's spines, a stolen pelt) add or subtract \textbf{dice}, not successes --- the 4+ line is bedrock.

\textbf{Zero rating.} You may usually still try: roll a single die, and the GM may charge an extra required success. Some checks are trained-only.

\textbf{Automatic information.} Some traits (Uln Icecunning, Brindel Brushcunning) grant automatic success on specific Perceive checks. That is the general pattern for ``automatic'' traits: inside the trait's territory, the dice concede. [PROPOSAL P6] Where an older trait references a check the sheet no longer lists (\textit{Track}, in the Squar's Living Myth), read it as \textbf{Survive} to follow a trail, or \textbf{Perceive} to notice one.

\subsection{The Twelve Skills}

\begin{description}
    \item[\textbf{Mobility}] --- \textbf{Athletics} (running, climbing, swimming, force in motion --- and the fuel gauge of Korr flight); \textbf{Knack} (hands and tools under pressure: locks, knots, sleight, jury-rigs); \textbf{Nimbleness} (balance, evasion, squeezing, grace --- the Tikka's and Dóm's home ground).
    \item[\textbf{Exploration}] --- \textbf{Investigate} (deliberate examination: scenes, documents, ledgers, lies already told); \textbf{Perceive} (the passive and the immediate: notice, spot, hear, read a room); \textbf{Survive} (weather, forage, trails, the country's rules --- and following what left sign).
    \item[\textbf{Knowledge}] --- \textbf{Study} (the learned disciplines: procedure, history, engineering, law); \textbf{Folklore} (what the road knows: customs, sign, stories that are load-bearing, which glyph means \textit{do not}); \textbf{Religion} (the Heavens, the churches, the rites --- and what the rites are \textit{for}).
    \item[\textbf{Social}] --- \textbf{Convince} (honest persuasion: appeal, negotiation, testimony); \textbf{Deceive} (the crafted impression: lies, disguise, the face-taker's second discipline); \textbf{Intimidate} (fear as an instrument --- the Nol's cultivated inheritance).
\end{description}

Class, species, and faction determine which of these your life trained; the sheet does not care how you came by them.

\medskip\hrule\medskip

\section{Alignments}

The five Alignments are the world's five natures, rated 0--10 like everything else, and arranged on your sheet as the pentacle they actually form.

\begin{center}
\begin{tabular}{ll}
\textbf{Heavens} & Compassion, law, morality \\
\textbf{Machine} & Cleverness, logic, technology \\
\textbf{Shadow} & Ambition, deception, ruthlessness \\
\textbf{Wild} & Growth, physicality, vigor \\
\textbf{Wyrd} & Chaos, curiosity, emotion \\
\end{tabular}
\end{center}

An Alignment rating is a \textbf{relationship}, not a virtue score. Three Shadow does not make you wicked; it means you and the Shadow have an understanding, and the GM is invited to remember it. Every character has some rating in all five, because every person alive is made of all five natures --- the ratings record which ones answer when you call.

Alignments do two things. [PROPOSAL P1 --- the pricing split below is proposed.]

\textbf{Augment (add dice).} When a Skill roll genuinely expresses an Alignment's nature --- Deceive carried by \textbf{Shadow}, an engineering Study leaning on \textbf{Machine}, Athletics surging with \textbf{Wild}, a plea to a mob that is honestly \textbf{Heavens} --- you may add dice equal to your Alignment rating to the roll. \textit{When applicable} is a real gate: the fiction must run through the plane's nature, and the GM arbitrates. Augmenting \textbf{invites} the plane: the GM \textit{may} attach a complication, in proportion to how hard you leaned --- three dice of Shadow in a lie is a whisper; eight is a partnership.

\textbf{Invoke (act outright).} When you set skill aside and act \textit{through} the plane itself, you do not roll. You receive \textbf{automatic successes equal to your rating} --- and it \textbf{always} costs you something, or causes another problem, as the GM determines. This is the game's central bargain, meant to be felt every time: the planes are reliable in a way training never is, and they are never free, because you are drawing on something with its own nature and its own accounts.

The cost should suit the plane, and a GM in need of the right complication can reach for its character:

\begin{itemize}
    \item \textbf{Heavens} collects in \textit{obligation}: the mercy you must now extend, the rule you must now keep, the witness you must now be. Its power arrives with its limits attached --- there are things it will not let you do \textit{next}.
    \item \textbf{Machine} collects in \textit{literalism}: the solution works exactly as specified and not one inch further; the time was real time; the logic that carried you now binds you --- this plan cannot be improvised away from.
    \item \textbf{Shadow} collects in \textit{trust}: someone saw, or suspects, or was spent; the ruthless path forecloses the honest one; and the Shadow's oldest surcharge --- what worked once now looks like the obvious tool for everything.
    \item \textbf{Wild} collects in \textit{excess}: the vigor overshoots --- collateral, noise, appetite, growth where you wanted none; the body remembers being let off the leash, and wants it again.
    \item \textbf{Wyrd} collects in \textit{strangeness}: the emotion escapes its handler; the curiosity opens a second door; something noticed the ripple in the dreaming, and tonight your dreams are shared property.
\end{itemize}

\textbf{Saving throws.} [PROPOSAL P2] When something targets you \textit{through} a plane's nature --- a bargain-hook (Shadow), a sovereign's terror (Wild), a Nightmare-thing's grip (Wyrd) --- roll dice equal to your rating in the targeted Alignment; 4+ succeeds; meet the effect's required successes. Saves are forced, so they carry no complication --- the price rule governs only what you choose. (This is the frame the Wyr Maker-Mark plugs into: a Wyr may save with their dominant Alignment's rating in place of their recessive one's.)

\medskip\hrule\medskip

\section{Momentum and the Flow of an Encounter}

\subsection{Sides}

When a situation turns into an \textbf{encounter}, the GM declares the \textbf{sides} --- usually two, the party and the opposition, though a three-way standoff or a battlefield with a hungry neutral in it plays exactly as it sounds. Sides \textbf{alternate activations}; the GM decides who surges first (ambushers, usually); one full cycle of all sides is a \textbf{round}, and the whole engagement is a \textbf{scene} for anything that recharges ``per scene.''

\subsection{The Surge}

Whenever a side activates:

\begin{enumerate}
    \item \textbf{Reload.} [PROPOSAL P4] Every character on the side draws from their holstered pile until their loaded hand holds 5.
    \item \textbf{Roll Momentum.} The side rolls \textbf{1d6} and gains that much Momentum. Characters and cards that grant \textbf{bonus Momentum} add it now --- the game's meaning of \textit{fast} is not ``acts first'' but ``makes the whole side richer.''
    \item \textbf{Spend.} The side spends its Momentum on actions, split among its members \textit{as they wish, in any order} --- one hero can burn the whole surge; six characters can act once each; anything between, freely interleaved.
    \item \textbf{Pass.} [PROPOSAL P3] When the Momentum is spent, or the side chooses to stop, the activation ends and unspent Momentum is lost. The top does not bank its spin.
\end{enumerate}

The d6 is the game's weather. A surge of 6 is a side moving like one animal; a surge of 1 is the terrible moment when only one thing can be done and the side must agree, fast, on what.

\subsection{Actions and Costs}

Unless stated otherwise, an action costs \textbf{1 Momentum}:

\begin{itemize}
    \item \textbf{Move} --- up to your \textbf{Move} rating in tiles.
    \item \textbf{Strike or Shoot} --- an attack (below).
    \item \textbf{Skill action} --- any other Skill roll: grapple, perceive, persuade, disable, aid.
    \item \textbf{Invoke} --- an Alignment action, automatic successes and price included. (Augmenting a Skill roll is not a separate action; it rides the roll it sharpens.)
    \item \textbf{Play a card} --- at the Momentum cost printed on it (usually 1).
\end{itemize}

\textbf{Long actions} cost more --- the GM prices them (hauling the gate shut: 2; rigging the charge properly: 3) --- and a side may pay a long action's cost across multiple activations if the fiction lets the work pause and resume. \textbf{Free actions} (dropping something, a shouted phrase, traits that say so --- Defiant Heart, Living Myth) cost nothing and ride alongside anything else.

\subsection{Combat: Strike, Shoot, Defend}

Attacks use the sheet's three combat stats, and the reminder printed on the sheet is the whole rule:

\begin{center}
\textbf{Roll Strike or Shoot dice. 4+ hits. Defend subtracts dice.}
\end{center}

To attack, roll dice equal to your \textbf{Strike} (melee) or \textbf{Shoot} (ranged), \textit{minus} the target's \textbf{Defend}. Every 4+ is a hit. [PROPOSAL P5 --- placeholder pending the combat chapter:] hits are damage; \textbf{Armor} reduces damage taken; damage comes off \textbf{HP}, and at 0 you are down and dying (a human, famously, lasts one round longer down there than anyone else). Traits and cards modify exactly where they say --- the Serekh's spines make even bare hands lethal and add $+1$ Defend against melee; the Gorûn's extra die joins any Strike built on raw strength. Weapons, conditions, and the rest of the arsenal belong to the combat chapter of the mechanical overhaul; this scaffold exists so the cards have something to hang on.

Note the design grain: Defend subtracting \textit{dice} means a hard-to-hit target makes your good pool small rather than your target number high --- fighting something with Defend 4 \textit{feels} like swinging at smoke --- and it means overwhelming an untouchable enemy is a Momentum problem (more attacks) or a card problem (bigger pools, ignored Defend), never a ``roll better'' problem. That is intended.

\medskip\hrule\medskip

\section{Powers: The Deck}

Your Powers are a deck of cards --- physical cards, kept beside the sheet --- and the deck's three piles use the only metaphor this setting would tolerate:

\begin{itemize}
    \item \textbf{Loaded} --- your hand, five cards: what's in the cylinder.
    \item \textbf{Holstered} --- your draw pile, face down: what you're carrying.
    \item \textbf{Expended} --- your discard: brass on the ground.
\end{itemize}

\textbf{At the start of an encounter}, load 5 from holstered. \textbf{When you play a card}, pay its Momentum cost and put it in expended --- unless it bears the \textbf{One-Shot} tag, in which case it is removed from play until the end of the encounter: some things fire once. \textbf{Whenever your side activates}, reload to 5; if holstered runs dry, \textbf{shuffle expended into holstered} and keep loading --- a long fight brings your tricks back around, in a new order, minus your One-Shots.

[PROPOSAL P9] \textbf{When the encounter ends}, everything --- expended and One-Shots alike --- returns to holstered. Between encounters, if a card's effect makes sense outside conflict (a Sawbones' craft, a key-mind's opened way), the GM may allow it narratively without touching the piles: the deck economy paces \textit{encounters}, not identity.

\textbf{Reading a card.} Every card carries its \textbf{name}, its \textbf{source} (class, species, or faction), its \textbf{level}, its \textbf{Momentum cost} (usually 1), its \textbf{effect}, and its \textbf{tags} (One-Shot is the first; the class sets bring more). What a card \textit{is} in the fiction depends on its source: a Captain's cards are her muster and her ledger arriving on time; an Arcanist's are craft carried in the chained book; a Dynast's are the blood doing what the blood does. The deck is not a spellbook bolted onto a character. The deck is the character, dealt out --- which is why the sheet doesn't duplicate it.

\subsection{Building the Deck}

Your \textbf{base deck is 20 cards}:

\begin{center}
\begin{tabular}{ll}
\textbf{10 cards} & from your \textbf{class pool} \\
\textbf{5 cards} & from your \textbf{species pool} \\
\textbf{5 cards} & from your \textbf{faction pool} \\
\end{tabular}
\end{center}

Each class publishes a \textbf{pool} --- everything it can ever teach --- organized by \textbf{card level}; at creation you choose from level-1 cards, taking duplicates where a pool permits them (the entry will say). A deck heavy in copies is consistent and narrow; a deck of singletons is versatile and unreliable; the shuffle will make you feel whichever you chose.

[PROPOSAL P7] \textbf{Leveling.} Each character level, \textbf{add 2 cards} from any of your pools at any card level up to your own, and \textbf{swap 1} existing card for another legal choice. The deck grows as you do --- and a bigger deck cycles slower, so taking every shiny card dilutes your best draws. Veterans prune. This is intended.

[PROPOSAL P8] \textbf{The unaffiliated.} A character with no faction --- rare among adventurers, and always a story --- replaces the 5 faction cards with 5 from the \textbf{Drifter pool}: the hard-won, unhoused competences the loose world teaches. GMs may also permit reaching into class or species pools for some of the five, at the price the Unaffiliated chapter describes: the nets stay open, and a Drifter whose deck starts looking like a faction's is a Drifter that faction will eventually come to see.

\textbf{Faction cards and leaving.} Faction cards are standing, training, and access as much as technique. If a character breaks with their faction, review the five with the GM: some are \textit{yours now} (the training took), some convert to Drifter equivalents, and one or two are simply gone --- the door that no longer opens. Severance should be legible on the character sheet. In this setting, it always is.

\medskip\hrule\medskip

\section{Identity, Deeds, and Possessions}

The sheet's remaining boxes are not bookkeeping; they are prompts the game will pull on.

\textbf{Demeanor, Drive, Reputation.} Three short lines: how you present, what moves you, and what the world says about you --- written by the player, revised in play. Demeanor is yours; Drive is yours; \textbf{Reputation is not} --- it is the word, carried mouth to mouth through the factors' web and the station taverns, and the GM owns its updates. When your Reputation and your Demeanor disagree, that is not an error on the sheet. That is a campaign.

\textbf{Nickname / Wanted-Poster Name.} The name the frontier uses when it is not being polite --- and, should it come to that, the name at the head of the writ. Players may leave it blank at creation. It rarely stays blank.

\textbf{Deeds.} [PROPOSAL P12 --- sketch.] Six slots for \textbf{specializations and signature abilities}, earned rather than bought: when a character accomplishes something the table agrees was \textit{a deed} --- held the door at Cold Hollow, crossed the Smoking Mother's ground in fire season, put a name back on a writ --- the GM may award a Deed: a named specialization granting an extra die in its exact circumstances, or a signature use of a card or trait. Deeds are written in the language of what happened (\textit{``Held at Redwater''}, not \textit{``+1 vs. crowds''}), because in this world --- the Barudrim have always been right about this --- what you have done is a stat.

\textbf{Coins and Funds.} [PROPOSAL P10] \textbf{Coins} are Rings in hand --- the Concord's neutral metal, good everywhere, heavy in the pocket. \textbf{Funds} are everything that isn't: banked credit, Company scrip, letters of draw --- money that is a \textit{relationship}, worth its issuer's health and reach, and subject to the spread. A party paid in Coins was paid. A party paid in Funds was promised, and the difference is a plot the Trade Companies chapter prices in detail.

\medskip\hrule\medskip

\section{An Exchange of Play}

\textit{The party --- Wren (Warden Shapeshifter), Ozanne (Lions Zealot), Pike (Guild Tinkerer, Coilwright-trained) --- is cornered at a relay post by three of the Rust Belt's organized things. Two sides. The GM gives the enemy the first surge; it spends 4 Momentum putting one thing on the roof and two at the door.}

\medskip

\textbf{GM:} Your side. Reload to five, everyone. Roll Momentum.

\textbf{Ozanne:} \textit{(rolls)} A 5. Okay --- Pike, the door?

\textbf{Pike:} One Momentum: I play \textbf{Overload Coil} --- cost 1 --- and dump my lantern's charge into the door bar. Card says 3 automatic damage to the first thing through. Expended.

\textbf{GM:} The bar starts humming. Four Momentum left.

\textbf{Wren:} One for me: I \textbf{shift} --- the hawk, a 1-cost species card --- and one more to \textbf{move}, out the smoke-hole, three tiles up. Eyes on the roof.

\textbf{GM:} Perceive --- it's dark and the thing is patient, so I want \textbf{two successes}.

\textbf{Wren:} Perceive 4, and I'll augment with \textbf{Wild} --- the hawk's eyes are exactly what Wild means. Wild 3, so seven dice. \textit{(rolls: 5, 4, 2, 6, 3, 4, 1)} Four successes --- two over.

\textbf{GM:} You leaned on the Wild lightly, so it just purrs --- this time. Two successes see it; your margin buys more: it hasn't seen \textit{you}, and it isn't watching the post at all. It's watching the \textit{other two}. Like a sergeant.

\textbf{Ozanne:} Two Momentum left. I spend one --- and not on a Skill. It's forty souls behind me and a door. I \textbf{invoke Heavens}. Rating 4: four automatic successes to hold this doorway, and the prayer is honest.

\textbf{GM:} Four successes: you will hold, this round, whatever comes. And the Heavens collect: the thing at the door is \textit{wounded already} --- something got to it before you --- and now that you have seen that, killing it has stopped being simple for you. One Momentum left?

\textbf{Ozanne:} Held. We keep it. ---

\textbf{GM:} It's lost --- spin doesn't bank. Enemy side surges...

\medskip

\textit{Note what the exchange shows: the side spent as one organism; the augment added dice and drew only a warning, while the invocation bought certainty and bought trouble; margin became fiction; and the unspent point died with the activation, as it should.}

\medskip\hrule\medskip

\section{Character Creation (Summary)}

[Arrays in step 4--5 are placeholders --- PROPOSAL P11.]

\begin{enumerate}
    \item \textbf{Choose your species} --- one of the seventeen peoples. Take its trait; note its 5-card pool.
    \item \textbf{Choose your class} --- your art and its 10-card allotment from the class pool.
    \item \textbf{Choose your faction} --- or, rarely, none. Take the 5 faction cards (or the Drifter set), and know who your people are.
    \item \textbf{Rate your Skills} --- distribute \textit{[X]} points across the twelve (creation cap \textit{[C]}; the 5+ and 8+ thresholds are meant to be \textit{earned}, and species traits that break caps say so).
    \item \textbf{Rate your Alignments} --- distribute \textit{[Y]} points across the five planes (creation cap \textit{[C]}). Every point is a relationship, and the GM is invited to remember them all.
    \item \textbf{Write your Identity} --- Demeanor, Drive, and a Reputation the GM co-signs. Leave the wanted-poster line blank if you dare.
    \item \textbf{Build your deck} --- 20 cards. Shuffle. Holster.
    \item \textbf{Derive} Strike, Shoot, Defend, Armor, Move, and HP from class and species per the overhaul's tables.
\end{enumerate}

\medskip\hrule\medskip

\section{Design Notes (for the table, not the fiction)}

\textbf{Momentum is a conversation.} The d6 surge forces the side to negotiate every activation --- who needs it, who can wait, whether the long action is worth starting --- and that negotiation \textit{is} the party, mechanically enacted. The variance is not noise; it is the frontier's honest answer about how much any given moment allows.

\textbf{The deck is identity under pressure.} Reloading to five every surge means you are always dangerous and never in full control of \textit{which} you is available. The shuffle asks, over and over: what do you do when the thing you wanted isn't loaded? Build the deck like a character, not like a solution.

\textbf{Two prices, one theology.} Skills are safe and uncertain. Augments are sharp and \textit{noticed}. Invocations are certain and \textit{owed}. That ladder --- from training, through borrowing, to conducting --- is the setting's cosmology expressed as a decision players make several times a session, and GMs should keep the prices specific, planar, and cumulative. A character who leans on the Shadow all campaign should end the campaign with the Shadow's opinion of them; and the same, more comfortingly and no less bindingly, for the Heavens.

\medskip\hrule\medskip
